%%%%%%%% PERSONAL DRAFT PROPOSAL -- NOT A TEAM DOCUMENT %%%%%%%%%%%%%%%%%%%%
\documentclass[11pt]{article}

\usepackage{microtype}
\usepackage[margin=1in]{geometry}
\usepackage{graphicx}
\usepackage{booktabs}
\usepackage{hyperref}
\usepackage{amsmath}
\usepackage{amssymb}
\usepackage{tabularx}
\usepackage{array}
\usepackage{tikz}
\usetikzlibrary{arrows.meta,positioning,shapes.geometric}
\usepackage{algorithm}
\usepackage{algorithmic}
\usepackage{fancyhdr}
\usepackage[capitalize,noabbrev]{cleveref}

\newcommand{\theHalgorithm}{\arabic{algorithm}}

% Evidence labels, same convention as the team's midterm report
\newcommand{\lab}[1]{{\footnotesize\textsf{\textbf{#1}}}}
\newcommand{\OBS}{\lab{OBSERVED}}
\newcommand{\DER}{\lab{DERIVED}}
\newcommand{\ASM}{\lab{ASSUMED}}
\newcommand{\DES}{\lab{DESIGNED}}
\newcommand{\DRAFT}{\lab{DRAFT}}

\pagestyle{fancy}
\fancyhf{}
\fancyhead[L]{\small\textsf{DRAFT -- personal proposal, not reviewed by the team}}
\fancyhead[R]{\small\textsf{MP-1}}
\fancyfoot[C]{\thepage}
\renewcommand{\headrulewidth}{0.4pt}

\title{\textbf{A Reserve-Paced Lyft Budget}\\
\large A draft extension to the Blue Jay Night Ride transfer-rule family}
\author{Miles Zhou\\
\small Team FourSight \quad$\cdot$\quad Department of Applied Mathematics and Statistics, Johns Hopkins University\\
\small EN.553.402\,/\,553.602 \quad$\cdot$\quad Instructor: Prof.\ Anthony J.\ Kearsley \quad$\cdot$\quad Fall 2026}
\date{Draft, \today}

\begin{document}
\maketitle

\begin{center}
\fbox{\parbox{0.92\textwidth}{\centering\small
\textbf{Status: \DRAFT.} This is a personal working proposal, drafted on the \texttt{miles-dev} branch. It is
\emph{not} a team decision, has not been reviewed, and nothing in it is implemented. It is logged as
\texttt{MP-1} in \texttt{docs/technical/decision\_log\_miles.md}, a personal log kept separate from the
team's \texttt{decision\_log.md} so that drafting this does not create merge conflicts. If the team agrees to
pursue it, it should be promoted into the team decision log under a real \texttt{D-}/\texttt{P-} id.}}
\end{center}

\begin{abstract}
The six Lyft transfer rules currently implemented (\texttt{src/sim/lyft.py}, R0--R3) each decide whether to
transfer a waiting request to Lyft using only that request's own state -- how long it has waited, or the
dispatcher's predicted pickup time. None of them know how much Lyft capacity has already been used, so none
can be budget-aware. This note proposes a layered extension -- a hard trip-count cap, a budget- and
demand-aware threshold, within-night pacing, demand-scaled nightly allocation, and a held-back monthly
reserve for abnormally bad nights -- that sits entirely downstream of the existing rules and requires no
change to the van dispatcher (C or G). It reuses evidence already measured by the team (T07, T09) rather than
introducing a trained predictive model, and it is explicitly scoped as additive and optional.
\end{abstract}

\section{Motivation}
\label{sec:motivation}

\texttt{docs/technical/scenario\_methodology.md} \S7b already names ``a Lyft budget'' as a candidate
upper-level constraint in the (unimplemented) bilevel policy framing. \DER\ This proposal is a concrete,
buildable version of that named-but-unbuilt thread, deliberately scoped to avoid the full bilevel machinery.

Three facts ground the design:
\begin{itemize}
  \item \textbf{No cross-request memory exists today.} \texttt{fires(rule, tau, t\_request, unassigned\_since,
    predicted\_pickup, lag\_s)} is a pure function of one request's state (Listing~\ref{lst:fires}). \OBS
  \item \textbf{No dollars are modelled} (D-45): Lyft is tracked only as operational trip counts. Any budget
    here must therefore be a \emph{trip-count} budget, not a dollar budget, to stay consistent with that
    decision. \DES
  \item \textbf{Demand is measurably uneven.} T09 gives the average night's hourly share of completed rides;
    T07 gives pooled weekday cancellation counts, a proxy for Lyft pressure (cancel count correlates with
    Lyft use at $r=0.88$). \DER\ Both are reused directly below -- no new model is fit.
\end{itemize}

\begin{algorithm}[h]
\caption{Current transfer check, per waiting request (\texttt{src/sim/lyft.py})}
\label{lst:fires}
\begin{algorithmic}
\REQUIRE rule, $\tau$, $t_{\text{request}}$, unassigned\_since, predicted\_pickup, $\ell$
\STATE $(\text{kind}, \theta) \leftarrow \text{RULES}[\text{rule}]$ \COMMENT{$\theta$ is a fixed constant today}
\IF{kind = elapsed}
  \RETURN $\tau - t_{\text{request}} \ge \theta$
\ELSIF{kind = predicted}
  \RETURN $\text{predicted\_pickup} - t_{\text{request}} > \theta$
\ENDIF
\end{algorithmic}
\end{algorithm}

\section{Proposed budget system}
\label{sec:system}

Five layers, each buildable on the last; none requires touching the dispatcher.

\subsection{Layer 1 -- hard cap}
A counter $\textit{lyft\_remaining}$, decremented on every booking. Once it reaches zero the rule behaves as
R0 for the rest of the night. Cheapest to build; a hard wall -- the last requests of a depleted night get no
help regardless of need.

\subsection{Layer 2 -- flexible threshold}
Replace the fixed $\theta$ in \texttt{RULES} with $\theta(\tau)$, tightening as budget runs low relative to
remaining demand:
\[
  \theta(\tau) \;=\; \theta_{\text{base}} \cdot f\!\left(\frac{\textit{lyft\_remaining}(\tau)}
  {\textit{remaining\_requests}(\tau)}\right), \qquad f \text{ increasing.}
\]
Same \texttt{fires()} call signature; one new rule kind alongside \texttt{elapsed}, \texttt{predicted},
\texttt{unassigned}, \texttt{van\_vs\_lyft}.

\subsection{Layer 3 -- demand-aware pacing}
Layer 2 \emph{is} a pacing throttle, made explicit:
\[
  \textit{target\_pace}(\tau) = \frac{\textit{lyft\_remaining}(\tau)}{\textit{remaining\_requests}(\tau)},
  \qquad
  \textit{actual\_pace}(\tau) = \frac{\text{trips used so far}}{\text{requests served so far}}.
\]
If $\textit{actual\_pace}(\tau) > \textit{target\_pace}(\tau)$, tighten; otherwise loosen. Structurally
identical to ad-spend budget pacing -- a known control pattern, not a new invention.

\subsection{Layer 4 -- demand-scaled nightly allocation}
A separate knob from Layers 2--3 (which flex the \emph{threshold} inside one fixed nightly budget): scale the
night's \emph{allowance} itself by that night's forecasted volume, in the same spirit as D-30's demand-scaled
fleet.

\subsection{Layer 5 -- spreading one monthly total across 30 nights}
In increasing sophistication: (i) a flat equal split, simple and wrong (wastes slack on quiet nights,
starves busy ones); (ii) a static split proportional to T07's own pooled weekday cancellation counts --
\emph{recommended first build}, no new estimation required; (iii) a rolling controller that tracks cumulative
spend against cumulative target across the month and adjusts the current night's allowance to compensate --
the natural follow-up once (ii) exists.

\section{Full formula}
\label{sec:formula}

\begin{align}
R &= r \cdot B, \qquad A = (1-r)\cdot B
  & \text{reserve } R \text{ and allocatable pool } A \text{ from monthly total } B \notag\\[2pt]
w_d &= \frac{\text{cancel}_d}{\sum_{d'} \text{cancel}_{d'}}, \qquad
  \textit{base}_n = A \cdot \frac{w_{d(n)}}{k_{d(n)}}
  & \text{weekday weight (T07) and night } n\text{'s base allowance} \notag\\[2pt]
\textit{remaining\_demand}(\tau) &= \sum_{h \ge \text{hour}(\tau)} S_h
  & \text{remaining share of the night, from T09's hourly curve } S_h \notag\\[2pt]
\theta(\tau) &= \theta_{\text{base}} \cdot f\!\left(\frac{\textit{actual\_pace}(\tau)}{\textit{target\_pace}(\tau)}\right)
  & \text{live threshold, replacing the constant in \texttt{RULES}} \notag
\end{align}

\textbf{Reserve trigger} (checked continuously):
\[
  \text{if}\quad \textit{actual\_cancel\_rate}(\tau) > \kappa \cdot \textit{expected\_cancel\_rate}(\tau)
  \quad\Rightarrow\quad R \leftarrow R - \Delta R, \qquad \textit{lyft\_remaining}(\tau) \leftarrow
  \textit{lyft\_remaining}(\tau) + \Delta R.
\]
$B,r,R,A$ are fixed once per month; $w_d, k_d, \textit{base}_n$ once per night, before it starts; $\theta(\tau)$
and the reserve check recompute live, every epoch.

\section{How it sits inside one epoch}
\label{sec:flow}

\begin{figure}[h]
\centering
\begin{tikzpicture}[
  node distance=0.55cm and 1.1cm,
  every node/.style={font=\small},
  box/.style={draw, rounded corners, align=center, minimum width=4.6cm, minimum height=0.8cm, fill=gray!6},
  dec/.style={draw, diamond, aspect=2.4, align=center, inner sep=2pt, fill=blue!5},
  term/.style={draw, align=center, minimum width=4.2cm, minimum height=0.8cm, fill=green!8},
  arr/.style={-{Latex[length=2mm]}}
]
\node[box] (solve) {Van solve (C or G) completes\\\footnotesize D-46: one Lyft pass follows};
\node[box, below=of solve] (theta) {Compute $\theta(\tau)$ live\\\footnotesize (pacing formula, \S\ref{sec:formula})};
\node[dec, below=of theta] (fires) {\texttt{fires()}\\true?};
\node[term, below left=0.5cm and -0.3cm of fires] (stay1) {Request stays\\with the van plan};
\node[dec, below right=0.5cm and -0.3cm of fires] (cap) {\textit{lyft\_remaining}\\$>0$?};
\node[term, below left=0.5cm and -0.3cm of cap] (book) {Book Lyft;\\decrement counter};
\node[dec, below right=0.5cm and -0.3cm of cap] (trig) {Reserve trigger\\fires?};
\node[term, below left=0.5cm and -0.9cm of trig] (refill) {Refill $\Delta R$;\\book Lyft};
\node[term, below right=0.5cm and -0.3cm of trig] (stay2) {Request stays\\with the van plan};

\draw[arr] (solve) -- (theta);
\draw[arr] (theta) -- (fires);
\draw[arr] (fires) -- node[left,xshift=-1pt]{no} (stay1);
\draw[arr] (fires) -- node[right,xshift=1pt]{yes} (cap);
\draw[arr] (cap) -- node[left,xshift=-1pt]{yes} (book);
\draw[arr] (cap) -- node[right,xshift=1pt]{no} (trig);
\draw[arr] (trig) -- node[left,xshift=-1pt]{yes} (refill);
\draw[arr] (trig) -- node[right,xshift=1pt]{no} (stay2);
\end{tikzpicture}
\caption{Where the budget machinery sits relative to the existing rule logic. The rule family itself
(R0--R3) is unchanged; the budget can only veto a ``yes'' (hard cap) or, via the reserve, turn a would-be
``no'' back into a ``yes'' on a night that is already running hot.}
\label{fig:flow}
\end{figure}

On a calm night $\theta(\tau)$ barely moves and the reserve never triggers -- the rules behave exactly as
they do today. The machinery only becomes visible on a night that is running hot, which is the point.

\section{Open questions and status}
\label{sec:open}

\begin{itemize}
  \item \textbf{Nothing here is built.} No code exists for any layer; \texttt{src/sim/lyft.py} is unchanged.
  \item $f(\cdot)$'s exact shape, $\kappa$, $\Delta R$, and $r$ are undetermined -- placeholders for a design
    discussion, not tuned values. \ASM
  \item T09's hourly curve is an \emph{average}-night template, not any real night's actual shape; pacing
    against it will not be exact on any given night. \DER
  \item T07's weekday cancellation weights are a proxy for Lyft pressure, not Lyft counts directly -- a
    defensible but indirect choice. \DES
  \item With one historical bad night in the April data (Apr 17, itself a supply-side disruption rather than
    a demand spike -- T03b), the reserve's size cannot be statistically justified from this data alone, only
    set by judgment -- the same limitation D-14 already encountered for event-aware forecasting.
\end{itemize}

\section*{References}
\begin{enumerate}
  \item \texttt{src/sim/lyft.py}; \texttt{src/sim/tradeoff.py} -- current rule implementation and the existing
    Pareto tracking of Lyft bookings per night.
  \item \texttt{docs/technical/decision\_log.md}, D-45 (no financial modelling), D-30 (demand-scaled fleet),
    D-14 (event-aware forecasting declined, $n=1$).
  \item \texttt{docs/technical/scenario\_methodology.md} \S7b -- Lyft budget named as an unimplemented
    upper-level constraint.
  \item \texttt{docs/evidence/insights/t07.md}, \texttt{t09.md} -- weekday cancellation counts and hourly
    completed-ride share, reused directly in \S\ref{sec:formula}.
  \item Grahn, L., Qian, X., \& Hendrickson, C. (2023). Optimizing first- and last-mile public transit
    services leveraging TNCs. \emph{Transportation}, 50(5) -- closest published analogue, on budget/cost
    -constrained TNC outsourcing for transit.
\end{enumerate}

\end{document}
